% Осень26!A39; нумерация внутри столбца без пустых строк.
% Источник: лекция 6, §2 (однородность); лекция 7, §4 и «Долги и упражнения»
\ticket{36}{Линейность интеграла Лебега}
\label{ticket:36}

\materialnote{Есть формулировка и однородность; доказательства аддитивности в лекциях 1--7 нет.}

\textbf{Формулировка.} Для суммируемых функций \(f,g\) и чисел \(\alpha,\beta\in\mathbb R\) функция \(\alpha f+\beta g\) суммируема и \[\int_X(\alpha f+\beta g)\,d\mu=\alpha\int_X f\,d\mu+\beta\int_X g\,d\mu.\]

\textbf{Статус в источниках.} Однородность доказана в лекции 6 среди простейших свойств интеграла. Аддитивность \(\int(f+g)=\int f+\int g\) используется в лекции 7, но её доказательство отложено. Полного доказательства линейности в предоставленных материалах нет.
